\chapter{Introduction and Project Vision}

\section{Project Overview}
GeoCrop is a satellite-based agricultural intelligence toolkit. It takes multispectral satellite
imagery of an agricultural region, identifies which land is cropland, classifies the crop growing
in each 30~m pixel, groups pixels into field regions with area, and (in FYP-2) estimates yield. The
results are published as map layers and statistics in a Web GIS portal, where users can switch
layers on and off, filter by crop, search fields, inspect values at any location and read
crop-area summaries.

The system is built on open data. Sentinel-2~\cite{sentinel2} provides the imagery, the USDA Cropland Data Layer
(CDL)~\cite{cdlsite,lin2022,hao2020} provides labelled crop maps for training and evaluation over eight U.S.
areas of interest (AOIs), and USDA NASS statistics will provide historical yield. GeoCrop is designed
as an academic, department-oriented prototype whose models, evaluation method and outputs are
transparent, and whose software architecture can carry many images and many models without
slowing down.

\section{Problem Statement}
The current crop-survey workflow relies on manual field inspection, crop identification, boundary
digitisation, area calculation and GIS data entry. This process is time-consuming, needs significant
human resources, is difficult to repeat frequently over large regions, and can introduce
inconsistencies in crop classification and field mapping. Satellite imagery observes agricultural
land frequently and with geographic reference, but turning those observations into crop maps,
crop types, field regions, acreage and expected yield still requires substantial processing and
expertise, and the results must be delivered in a form that planners can use interactively over
large areas.

There is therefore a need for an integrated system that learns crop classes from labelled
geospatial data, analyses new satellite imagery, distinguishes cropland from other land, classifies
supported crop types, derives field regions and area, predicts yield from suitable spatio-temporal
features, and presents the results through a responsive Web GIS portal that remains fast as the
number of images and model layers grows.

\section{Problem Context and Motivation}
Agricultural departments, field survey teams, GIS analysts and planners need timely information
about what is planted, where, and over how much area. In the current workflow (physical survey,
crop and non-crop identification, crop-type identification, boundary drawing, area and coordinate
calculation, GIS entry, separate yield assessment) every step is manual. This is slow, needs many
people, is hard to scale across large areas, and can introduce errors.

Open Earth-observation missions now give repeat multispectral coverage, and machine learning has
shown that crop classes can be learned from spectral and temporal patterns~\cite{khan2023,
muqaddas2024,tran2022}. Research in Pakistan confirms that Sentinel-2 can identify crops in
smallholder settings, but still depended on thousands of manually surveyed field boundaries
\cite{khan2023}. In the United States, the CDL offers large labelled coverage that has been used to
train classifiers~\cite{hao2020,wang2020}. GeoCrop uses this labelled coverage to develop and
evaluate the methodology before adaptation to local data.

The motivation is also practical. A crop map is only useful if people can work with it. The
proposal-stage prototype showed that a single AOI's crop-type result, shipped as vector polygons,
already made the map slow to load and zoom. A system meant for many images and several models
(cropland, crop type, boundaries, yield) has to be designed for performance from the start. Better
agricultural monitoring supports informed crop planning and food-security decisions (SDG~2,
Zero Hunger).

\section{Project Vision}
GeoCrop aims to give agricultural organisations a single toolkit that turns new satellite imagery
into trustworthy, map-ready crop intelligence in minutes: where the cropland is, what is growing,
how large each crop area is, and what yield to expect. Results should be honest about their
accuracy, reproducible, and fast to explore on a map, and new models should plug into the same
platform as research progresses.

\section{Project Objectives}\label{sec:goals}
After the proposal defence the panel asked the team to clarify its goals (P1,
Section~\ref{sec:feedback}). The goals are now stated as follows. GeoCrop's primary goal is a
\emph{light, fast and integrated} platform: one that turns an uploaded Sentinel-2 image into crop
layers on a map without heavy hardware, and that stays usable as images and models grow. Its
research goal is crop-type models that are accurate under a fair test and cheap enough to serve.
Three terms from the proposal are defined precisely:
\begin{itemize}
  \item \textbf{Crop mapping}: per-pixel cropland and crop-type maps at 30~m for supported AOIs,
        scored against the CDL on spatially separated test blocks.
  \item \textbf{GIS integration}: every result is a georeferenced layer in one catalogue, viewable,
        filterable and comparable on the same map as the imagery and other layers.
  \item \textbf{Real-time estimation}: \emph{near} real time on demand. An AOI-sized image is analysed
        and published in about a minute (57~s measured), and a map screen loads in well under a
        second. Continuous streaming of new imagery is not a goal.
\end{itemize}
Each objective below has a measurable target that is checked in Chapter~7.
\begin{enumerate}[label=\textbf{O\arabic*.}, leftmargin=1.1cm]
  \item Build a spatially aligned Sentinel-2 and CDL training dataset for at least two AOIs, with
        a reproducible preparation pipeline.
  \item Develop and evaluate cropland and crop-type models using a spatially separated test set,
        reporting accuracy, macro-F1 and per-class scores with confidence intervals, and compare
        them with published results.
  \item Derive field regions and area from model output and publish them as GIS layers with
        geographic coordinates; add a dedicated field-boundary model in FYP-2.
  \item Develop a yield-estimation module evaluated with $R^2$, RMSE and MAE against USDA NASS
        records at the spatial scale for which ground truth exists.
  \item Deliver a Web GIS portal that shows one screen of results using under 1~MB of data, returns
        cold map tiles in under 300~ms and cached tiles in under 50~ms, and keeps several layers
        visible at once.
  \item Provide a processing pipeline in which an uploaded image is validated, analysed by the
        selected models and published to the map automatically, and in which a new model can be
        added without changing the API, database schema or user interface.
\end{enumerate}

\section{Project Scope}
\textbf{Included.}
\begin{itemize}
  \item Satellite data preparation: Sentinel-2 composites, band selection, reprojection to the
        CDL grid, vegetation indices and monthly time series.
  \item CDL-based label preparation for the supplied AOIs, including rare-class grouping.
  \item Cropland/non-cropland and multi-class crop-type mapping with model comparison.
  \item Field-region derivation, area statistics and vector map layers; a field-boundary model in
        FYP-2.
  \item Yield estimation using satellite features and USDA NASS history (FYP-2).
  \item A Web GIS platform: upload and processing jobs, layer catalogue, tile-based map, dashboard,
        model catalogue, system status, and user accounts.
\end{itemize}
\textbf{Excluded.} Building or operating satellites; legal or cadastral ownership mapping; crop
disease diagnosis; real-time farm machinery integration; production-scale national deployment
with service-level guarantees; and claims of ground-truth accuracy beyond agreement with the
reference labels.

\section{Project Modules}
\begin{table}[htbp]
\centering\small
\caption{GeoCrop modules and responsibilities}
\label{tab:modules}
\begin{tabularx}{\textwidth}{L{4.2cm} Y L{2.2cm}}
\toprule
\textbf{Module} & \textbf{Responsibility} & \textbf{Status}\\
\midrule
M1 Data acquisition and preprocessing & Export Sentinel-2 composites, align them to the CDL grid (EPSG:5070, 30~m), build monthly time series and indices, convert results to Cloud Optimized GeoTIFF. & Working for AOI1, AOI2\\
M2 Crop mapping & Cropland/non-cropland and crop-type models; evaluation and model selection. & Baseline in platform; research pipeline selected\\
M3 Field regions, boundaries and area & Polygonise model output, compute area, build vector tiles; field-boundary model later. & Field regions working; model in FYP-2\\
M4 Yield estimation & Regression on satellite features and NASS yield history. & FYP-2 (registered placeholder)\\
M5 Processing pipeline & Upload, validation, preprocessing, background model runs, progress tracking. & Working\\
M6 Model registry & Plug-in contract, manifests, dependency resolution, model catalogue. & Working\\
M7 Catalogue, storage and statistics & PostGIS catalogue of areas, runs, layers, statistics and fields; S3 object storage. & Working\\
M8 Web GIS portal & Home, Live Map, Dashboard, Analyze Imagery, Models, System Status, accounts. & Working\\
\bottomrule
\end{tabularx}
\end{table}

\section{Stakeholders and User Classes}
\begin{table}[htbp]
\centering\small
\caption{Stakeholders and user classes}
\label{tab:stakeholders}
\begin{tabularx}{\textwidth}{L{4.2cm} Y}
\toprule
\textbf{Stakeholder / user class} & \textbf{Characteristics and needs}\\
\midrule
Agriculture Department & Needs regional crop area and expected production for planning; values accuracy and transparency.\\
Field survey teams & Today collect crop and boundary data manually; benefit from fewer field visits and targeted checks.\\
GIS analysts & Technical users; need georeferenced layers, coordinates, area and exportable data.\\
Agricultural planners and decision makers & Non-technical; need clear maps, legends and summary statistics.\\
Farmers & Indirect beneficiaries of better planning and advisory services.\\
Industry / startup partner & External collaborator interested in a scalable, light system.\\
Platform developers and ML engineers & Add and update models, monitor the system, run evaluations.\\
\bottomrule
\end{tabularx}
\end{table}

\section{Success Criteria}
\begin{itemize}
  \item \textbf{Crop mapping:} crop-type accuracy and macro-F1 reported on spatially separated test
        blocks with 95\% intervals; results comparable in order of magnitude to published CDL-based
        studies for similar regions (for example 94\% South Dakota and 83\% California
        in~\cite{tran2022}).
  \item \textbf{Boundaries and area:} spatial overlap (IoU) and area error against reference data,
        and acreage within a stated tolerance of USDA statistics (FYP-2).
  \item \textbf{Yield:} $R^2$, RMSE and MAE on held-out years or counties (FYP-2).
  \item \textbf{Performance:} first visible result under 1~s, under 1~MB per screen, cold tiles
        under 300~ms, cached tiles under 50~ms, eight or more layers visible.
  \item \textbf{Pipeline:} an uploaded image reaches the map without manual steps, with readable
        errors for invalid input; a new model is added by folder and manifest only.
  \item \textbf{Lightweight:} the full workflow runs on one laptop-class host; the selected crop-type
        model is CPU-only and small; peak memory per stage is measured.
  \item \textbf{Usability:} non-technical users can find, filter and read results using legends and
        plain-language labels; measured with task completion and a questionnaire score in a user
        study.
  \item \textbf{Integration:} several layers from different models and images can be used together on
        one map without loading whole datasets.
\end{itemize}

\section{Constraints and Assumptions}
\begin{itemize}
  \item CDL labels are a classified product, not field survey data; scores measure agreement with
        the CDL~\cite{lin2022}.
  \item Sentinel-2 at 10--20~m, resampled to the 30~m CDL grid, limits boundary precision for small
        fields.
  \item Data for 2022 is used so that imagery and labels refer to the same season.
  \item Yield ground truth (USDA NASS) is available at county level, which bounds the spatial scale
        of yield evaluation.
  \item Compute is limited to student laptops and free cloud notebooks, and the panel named a
        lightweight system as a main concern, so CPU-only inference is preferred and resource use
        is measured for every stage (Section~\ref{sec:feasibility}).
  \item Uploaded imagery is assumed to be a georeferenced Sentinel-2 surface-reflectance stack with
        the documented bands; other formats are rejected with a clear message.
  \item The platform runs on Docker Compose on a single host for development and demonstration.
\end{itemize}

\section{Ethical, Legal and DEI Considerations}
\begin{itemize}
  \item \textbf{Honest reporting:} predictions are never presented as ground truth. The portal
        shows model accuracy, its evaluation protocol and limitations; placeholder models are
        labelled ``coming soon'' and never produce data.
  \item \textbf{Fairness across classes and regions:} per-class scores and per-area results are
        reported so that unequal performance (for example on small or rare crops) is visible.
  \item \textbf{Licensing:} Sentinel-2 data (Copernicus open licence), CDL and NASS data (U.S.
        public domain), OpenStreetMap basemap (ODbL attribution shown) and open-source libraries
        are used within their terms and acknowledged.
  \item \textbf{Privacy:} uploaded imagery and results belong to the uploading account; passwords
        are hashed; startup-provided data will be protected and not published.
  \item \textbf{Accessibility and inclusion:} the interface uses clear legends, plain-language
        labels, keyboard-accessible controls and readable colour palettes so that technical and
        non-technical staff can use it.
  \item \textbf{Disclosure:} external datasets, models, libraries and AI-assisted work are
        disclosed (Sections~\ref{sec:tools-assisted} and~\ref{sec:external}).
\end{itemize}
